\documentclass[11pt]{article}

\usepackage[margin=1in]{geometry}
\usepackage{amsmath,amssymb,amsthm}
\usepackage{float}      % provides the [H] float specifier (force-here placement)
\usepackage{enumitem}
\usepackage{xcolor}
\usepackage[colorlinks=true,linkcolor=blue!55!black,urlcolor=blue!55!black,citecolor=blue!55!black]{hyperref}

\newtheorem{theorem}{Theorem}[section]
\newtheorem{lemma}[theorem]{Lemma}
\newtheorem{proposition}[theorem]{Proposition}
\newtheorem{corollary}[theorem]{Corollary}
\newtheorem{definition}[theorem]{Definition}
\newtheorem{remark}[theorem]{Remark}

\newcommand{\HH}{\mathbb{H}}
\newcommand{\ZZ}{\mathbb{Z}}
\newcommand{\CC}{\mathbb{C}}
\newcommand{\RR}{\mathbb{R}}
\newcommand{\QQ}{\mathbb{Q}}
\newcommand{\NN}{\mathbb{N}}
\newcommand{\SLZ}{\mathrm{SL}_2(\ZZ)}
\newcommand{\Dhat}{\widehat{\Delta}}
\newcommand{\tk}{\widetilde{k}_T}
\DeclareMathOperator{\Res}{Res}
\DeclareMathOperator{\im}{im}
\DeclareMathOperator{\Li}{Li}

\title{$L$-functions for all meromorphic $\SLZ$ modular forms,\\
via manifestly finite contours}
\author{David A.\ McGady}
\date{}

\begin{document}
\maketitle

\begin{abstract}
We give a single manifestly finite, regulator-free functional realizing in one object $L$-functions
quasi-periods and period polynomials of meromorphic modular forms on $\SLZ$, with generic 
poles throughout the upper half-plane $\HH$ and at the cusp of $\SLZ$. Additionally, our construction 
is preserves holomorphicity. It is a two-segment contour integral, 
$L^*(f,s) = L^*(f, s; \gamma^T_{\tau_0}) + L^*(f, s; \gamma^S_{\tau_0})$, where the 
$\gamma^{S}_{\tau_0}$ is a path in $\HH$ from $S(\tau_0)$ to $\tau_0$, and the same for $\gamma^T_{\tau_0}$. 
For weakly holomorphic forms our L-function is independent of $\tau_0$ and of the path, and reproduces 
classical L-functions for holomorphic modular forms, reproduces the L-functions of weak modular forms 
in~\cite{BFK2011}, which previously needed non-holomorphic regularization to be defined, and reproduces 
Brown's quasi-periods~\cite{Brown2017}. For properly meromorphic modular forms with poles inside the 
interior of $\HH$, our L-function depends on the homotopy classes, defined in terms of pole-crossings. 
There is a natural choice of trivial homotopy class, which renders our L-functions canonical when the 
modular form has a pole in the interior of the fundamental domain, generalizing~\cite{1806}.
\end{abstract}

\section{Introduction and preliminary concepts}

Our primary goal in this note is to put forth a unifying generalization of the L-function 
associated to a given modular form. Our L-function naturally encompasses the classical 
case of L-functions associated to holomorphic modular forms~\cite{01-Apostol}, and 
seamlessly unifies Ref.~\cite{BFK2011}'s extension to L-functions of modular forms with 
poles at the cusp with Ref.~\cite{1806}'s extension to L-functions of modular forms with 
poles at finite locations, exclusively away from cusps. Secondarily, we show that our 
L-functions reproduce the numerical results for quasi-periods associated to weak cusp forms~\cite{Brown2017}/

\subsection{Modular forms}
Write ${\rm Re}(z)$ and ${\rm Im}(z)$ for the real and imaginary parts of a complex number 
$z\in \CC$. Let $\HH:=\{\tau\in\CC\mid {\rm Im}(\tau)>0\}$ denote the upper half-plane. $\SLZ$ 
is the group freely generated by $T := \left( \begin{smallmatrix} 1 & 1 \\ 0 & 1 \end{smallmatrix} \right)$ 
and $S := \left( \begin{smallmatrix} 0 & -1 \\ 1 & 0 \end{smallmatrix} \right)$. 

A {\em holomorphic modular form} of weight $k$ for $\SLZ$ is a 
holomorphic function $f:\HH\to \CC$ such that $f(\tau+1)=f(-\frac1\tau)\tau^{-k}=f(\tau)$ for 
all $\tau\in \HH$, and $f(\tau)={\cal~O}(1)$ as ${\rm Im}(\tau)\to \infty$, and which has no 
poles at finite locations within $\HH$. By virtue of the periodicity $f(\tau) = f(\tau+1)$ and 
holomorphicity, each holomorphic modular form has a q-series decomposition,
\begin{equation}
f(\tau) = \sum_n c_f(n) e^{2 \pi i n \tau} = \sum_n c_f(n) q^n~~,~~q:= e^{2 \pi i \tau}~~.
\end{equation}
In order for $f$ to be bounded as $\tau \to i \infty$, we conclude that $c_f(n) = 0$ 
for all $n < 0$, i.e. that $f = \sum_{n \geq 0} c_f(n) q^n$; for such modular forms $f$ 
we write $f \in M_k$. If a modular form $f$ additionally has a vanishing constant term, 
i.e. if $c_f(n) = 0$, then $f$ is a {\em cusp form}, 
and we write $f \in S_k$. Clearly, $S_k \subset M_k$. Further, $M_k$ is a  
finite-dimensional vector-space, as is $S_k$. Holomorphic $\SLZ$ modular forms must 
have even integer weights, and $k \geq 4$~\cite{01-Apostol}. 

A product of modular forms of weights $k_1$ and $k_2$ is a modular form of weight 
$k_1+k_2$, so $M_*:=\bigoplus_k  M_k$ is naturally a (graded) ring. This graded ring is 
freely generated by the two modular forms $E_4 \in M_4$ and $E_6 \in M_6$, the $k = 4, 6$
specializations of the {\em weight $k$ Eisenstein series} $E_k$, which may be defined as follows:
\begin{align}
E_k := \sum_{(m,n)=1} \frac{1}{(m+n\tau)^k} = 1 + \frac{1}{\zeta(1-k)}\sum_{n = 1}^{\infty} \sigma_{k-1}(n) q^n~,
\end{align}
where $\sum_{(m,n) = 1}$ denotes summing over all relatively prime $(m,n) \in \ZZ^2$,  $\zeta(s)$ is the 
Riemann zeta-function analytically continued to the entire $s$ plane, and 
$\sigma_m(n) := \sum_{d|n} d^m$ is the Divisor function where $d|n$ means $n/d \in \NN$. 
In short, $M_* = \CC[E_4, E_6]$. We now define the {\em modular discriminant}:% $\Delta \in S_{12}$:
\begin{equation}
\Delta := \frac{E_4^3-E_6^2}{1728} = q + O(q^2) \in S_{12}~.
\end{equation}

Holomorphicity is most canonically relaxed via allowing a modular form to have a pole at the 
cusp of $\SLZ$, i.e. for a weight $k$ modular form for $\SLZ$ to have negative powers of $q$ 
in its q-series expansion. In this case, $f$ diverges as $\tau \to i \infty$ as $f \sim e^{-2 \pi N \tau}$, 
where $-N$ is the largest negative number such that $c_f(n)$ is non-zero. Such a modular form 
$f$ is called a {\em weak modular form}, and we write $f \in M_k^!$. If a weak modular form $f$ 
has a vanishing coefficient of $q^0$, i.e. if $c_f(0) = 0$, it is called a {\em weak cusp form} and 
we write $f \in S_k^!$. Clearly, $S_k^! \subset M_k^!$, $S_k \subset S_k^!$ and $M_k \subset M_k^!$. 
Two important weak cusp forms in our paper are as follows: 
\begin{align}
J :=&~ \frac{E_4^3}{\Delta} - 744 = \frac{1}{q} + O(q) \in S_0^!~,~\\
\Dhat =&~ (J^2 + a_1 J + a_0) \Delta := \frac{1}{q} + O(q^2) \in S_{12}^!~,~
\end{align}
where $a_0, a_1$ are integers which exist to ensure the q-series gap that defines $\Dhat$~\cite{DukeJenkins}.

In this work, we further relax from even weakly holomorphic modular forms, and 
consider {\em meromorphic modular forms} of weight $k$ which are quotients: $f=\frac{f_1}{f_2}$ 
where $f_i\in M_{k_i}$, the denominator $f_2$ is not identically zero, and $k=k_1-k_2$. We 
write $F_k$ for the vector space of meromorphic modular forms of weight $k$, and set 
$F_*:=\bigoplus_k F_k$. Then $F_k$ vanishes unless $k$ is an even integer, and 
$F_*$ may be regarded as a ``graded ring of fractions'' of $M_*$. 

\subsection{L-functions for holomorphic modular forms and periods for cusp forms}

For any $f \in S_k$, one may define the associated L-function:
\begin{align}
L(f,s) := \frac{(2 \pi)^s}{\Gamma(s)} L^*(f,s) := \frac{(2 \pi)^s}{\Gamma(s)} \int_0^{\infty} f(it) t^{s-1} dt = \sum_{n \geq 1} \frac{c_f(n)}{n^s}~.
\end{align}
A similar definition applies for $f \in M_k$. %When $f \in M_k$, we define $L^*(f,s)|_{\rm classical} := 
%\int_0^{\infty} (f(it)-c_f(0)) t^{s-1} dt$. With the $-c_f(0)$-subtraction, then the L-integral is OK. However,  
%if $f$ has a constant term, then the constant term in the commuted q-series/integral diverges:  
%$\int_0^\infty t^{s-1} dt$ is a pure divergence. It is important e.g. in BFK and in my 1806-paper that  
%the constant term IS thus subtracted/out. Further, in doing so, we then get the dividend of the  
%$-c_f(0)/s --i^k c_f(0)/(k-s)$-terms in the L-integral, which I needed in order to establish my sum-rules  
%in 1806. Yet, in our manifestly finite construction, this projection is both not needed and never even  
%shows-up. We just naturally get the right answer! VERY pretty
Similarly, one may define a period polynomial for $f \in S_k$ in terms of this L-function,
\begin{align}
r_f(X, Y) := (2 \pi i)^{k-1} \sum_{\ell = 0}^{k-2} \binom{n}{\ell}\frac{L^*(f,\ell+1)}{i^{\ell+1}}\,X^{n-\ell}Y^\ell~.
\end{align}
This polynomial, $r_f(X, Y) \in V_{k-2} := \CC[X, Y]_{k-2}$, has a natural action under $\SLZ$, 
$(P|_{\gamma})(X, Y) = P(a X + b Y, c X + d Y)$ for $\gamma = 
\left( \begin{smallmatrix} a & b \\ c & d \end{smallmatrix} \right) \in \SLZ$. This period polynomial 
can be decomposed into a basis of highly constrained polynomials, which are annihilated by $|(1+T)$ 
and $|(1 + U + U^2)$\cite{KZ1984}, where $U := ST$ generates an order-3 cyclic subgroup of $\SLZ$. 
We define the \emph{periods} of a cusp form to be the coefficients that localize $r_f$ within this highly 
constrained sub-space of $V_{k-2}$.

However, as classically computed, these L-functions and associated period polynomials 
(when defined in terms of the special values of the associated L-function) are not defined when $f$ 
has poles either at the cusp of $\SLZ$ or within the interior of $\HH$. As stated above, there 
are extant, but partial and non-overlapping, extensions of the definition of $L(f,s)$ to $f \in M_k^!$ or 
$f \in F_k$. In this note, we provide a construction that unifies these disparate extensions.

\section{A finite L-integral}\label{sec:Lint}

\begin{definition}[$S$- and $T$-segments]\label{def:segments}
Let $\tau_0\in\HH$ and $f\in F_k$, let
$S=\left(\begin{smallmatrix}0&-1\\1&0\end{smallmatrix}\right)$,
$T=\left(\begin{smallmatrix}1&1\\0&1\end{smallmatrix}\right)$ denote 
the generators of $\Gamma=\SLZ$, and let $f^{-1}(\infty)$ be the locus 
of all poles of $f \in \HH$ (inclusive of the cusp of $\SLZ$).  
The \emph{$S$- and $T$-segments at $\tau_0$} are smooth paths in
$\HH\setminus\Gamma\!\cdot\!f^{-1}(\infty)$,
\begin{equation}
\gamma^S_{\tau_0}:\ -1/\tau_0\ \longrightarrow\ \tau_0,\qquad
\gamma^T_{\tau_0}:\ \tau_0-1\ \longrightarrow\ \tau_0,
\end{equation}
running from $S^{-1}\tau_0=-1/\tau_0$ and $T^{-1}\tau_0=\tau_0-1$ to $\tau_0$.
\end{definition}

Avoiding poles is automatic for $f \in M_k^!$, as long as the base-point $\tau_0$, 
is not at a cusp of $\SLZ$, i.e. there is only the trivial homotopy class. 
When $f \in F_k$, and $f$ has poles at finite and non-zero imaginary heights, 
integration contours fall into nontrivial homotopy classes.

\begin{definition}[Hurwitz $T$-kernel]\label{def:kernel}
Let $\zeta(s,x)=\sum_{m\ge0}(m+x)^{-s}$ as a function 
$\zeta: \CC_s \times (\CC\setminus\NN)_x \to \CC$ denote the 
Hurwitz zeta function (meromorphic in $s$, holomorphic in $x\in\HH$). 
We define \emph{Hurwitz $T$-kernel} as
\begin{equation}
\tk(\tau,s):=\zeta(1-s,\tau+1)-e^{i\pi(s-1)}\,\zeta\big(1-(k-s),\tau+1\big).
\end{equation}
\end{definition}

\begin{definition}[$L$-integral]\label{def:Lint}
Let $f\in F_k$, $\tau_0\in\HH$, and $s\in\CC$. Further, let $\gamma^S_{\tau_0}$,
$\gamma^T_{\tau_0}$ be the $S$- and $T$-segments at $\tau_0$, as in
Definition~\ref{def:segments}, and let $\tk$ be as defined in 
Definition~\ref{def:kernel}. The \emph{$L$-integral} of $f$ is
\begin{equation}
L^*(f,s;\gamma^S_{\tau_0},\gamma^T_{\tau_0})
:=e^{-\pi i s/2}\Big[\int_{\gamma^S_{\tau_0}}\! f(\tau)\,\tau^{s-1}\,d\tau
+\int_{\gamma^T_{\tau_0}}\! f(\tau)\,\tk(\tau,s)\,d\tau\Big].
\end{equation}
\end{definition}

We use this \emph{L-integral} to define a generalized 
\emph{L-function}:

\begin{definition}[$L$-function]\label{def:Lfun}
Let $L^*(f,s;\gamma^S_{\tau_0}, \gamma^T_{\tau_0})$ be as defined in~\ref{def:Lint}. 
We define the generalized \emph{$L$-function} of $f\in F_k$ to be,
\begin{equation}
L(f,s;\gamma^S_{\tau_0},\gamma^T_{\tau_0})
:=\frac{(2\pi)^s}{\Gamma(s)}\,L^*(f,s;\gamma^S_{\tau_0},\gamma^T_{\tau_0}).
\end{equation}
\end{definition}

Having defined the L-integral and L-function, we prove they are manifestly finite for $f \in F_k$: 

\begin{lemma}[Finiteness]\label{lem:finite}
For $f\in F_k$, both integrals in Definition~\ref{def:Lint} are finite for every
$s\in\CC$, every $\tau_0\in\HH$, and every choice of paths
$\gamma^S_{\tau_0},\gamma^T_{\tau_0}$.
\end{lemma}

\begin{proof}
Both segments are compact and lie in $\HH\setminus\Gamma\!\cdot\!f^{-1}(\infty)$, so $f$ is
continuous on them; $\tau^{s-1}$ is continuous on $\HH$, and $\tk(\tau,s)$ is continuous there
since $\zeta(\nu,\tau+1)$ is holomorphic in $\tau\in\HH$ (Definition~\ref{def:kernel}). A
continuous integrand on a compact path is integrable, for every $s\in\CC$.
\end{proof}

We now turn to path-dependence and $\tau_0$-independence of 
the L-integral and the L-function of $f \in F_k$. 
%For $f \in F_k$ such that $L^*(f,s; \gamma^S_{\tau_0}, \gamma^T_{\tau_0})$ 
%is \emph{both} $\tau_0$-independant and path-independent, and we simply 
%write $L^*(f, s: \gamma^S_{\tau_0}, \gamma^T_{\tau_0}) = L^*(f,s)$; similarly for $L(f,s)$. 
This holds for all $f \in M_k^!$:

\begin{proposition}[Basepoint and path independence]\label{prop:indep}
For $f\in M^!_k$ the value $L^*(f,s;\gamma^S_{\tau_0},\gamma^T_{\tau_0})$ is independent
of $\tau_0\in\HH$ and of the paths: with no interior poles, $\HH$ is simply connected,
so all admissible contours are homotopic. We write $L^*(f,s)$ for the common value and
$L(f,s):=\frac{(2\pi)^s}{\Gamma(s)}L^*(f,s)$.
\end{proposition}

\begin{proof}
With no interior poles the integrands $f(\tau)\tau^{s-1}$ and $f(\tau)\tk(\tau,s)$ are
holomorphic on $\HH$, which is simply connected; by Cauchy's theorem each integral is unchanged
under any deformation fixing its endpoints, hence depends on those endpoints alone. This is
path-independence.

For independence of $\tau_0$ we differentiate the two integrals and watch them cancel. The
$T$-segment runs from $T^{-1}\tau_0=\tau_0-1$ to $\tau_0$ --- two endpoints a unit apart,
moving together with $\tau_0$ --- so by Leibniz's rule and the $T$-periodicity $f(\tau_0-1)=f(\tau_0)$,
\begin{equation}
\partial_{\tau_0}\!\!\int_{\gamma^T}\!\! f\,\tk
=f(\tau_0)\,\tk(\tau_0,s)-f(\tau_0-1)\,\tk(\tau_0-1,s)
=f(\tau_0)\,\delta_T\tk(\tau_0-1,s),
\end{equation}
where $\delta_T\tk(\tau,s):=\tk(\tau+1,s)-\tk(\tau,s)$. This is the whole role of $\delta_T$: the
unit gap between the endpoints turns the $\partial_{\tau_0}$ boundary term into a unit shift of
the kernel. That shift is elementary --- the kernel was built for it --- since
$\zeta(\nu,x+1)-\zeta(\nu,x)=-x^{-\nu}$ gives, on both terms at $x=\tau+1$,
\begin{align}
\delta_T\tk(\tau,s)&=-(\tau+1)^{s-1}+e^{i\pi(s-1)}(\tau+1)^{k-1-s}
\\
\implies \partial_{\tau_0}\!\!\int_{\gamma^T}\!\! f\,\tk&=f(\tau_0)\big[-\tau_0^{s-1}+e^{i\pi(s-1)}\tau_0^{k-1-s}\big].
\end{align}
The $S$-segment runs $-1/\tau_0\to\tau_0$; its lower endpoint moves with
$\partial_{\tau_0}(-1/\tau_0)=\tau_0^{-2}$, so by Leibniz, the modularity
$f(-1/\tau_0)=\tau_0^k f(\tau_0)$, and the phase $(-1/\tau_0)^{s-1}=e^{i\pi(s-1)}\tau_0^{1-s}$,
\begin{equation}
\partial_{\tau_0}\!\!\int_{\gamma^S}\!\! f\,\tau^{s-1}
=f(\tau_0)\tau_0^{s-1}-f(-1/\tau_0)\,(-1/\tau_0)^{s-1}\,\tau_0^{-2}
=f(\tau_0)\big[\tau_0^{s-1}-e^{i\pi(s-1)}\tau_0^{k-1-s}\big].
\end{equation}
The two brackets are negatives of each other, so
$\partial_{\tau_0}L^*=e^{-\pi i s/2}\big(\partial_{\tau_0}\!\int_{\gamma^S}+\partial_{\tau_0}\!\int_{\gamma^T}\big)=0$.
\end{proof}

\begin{remark}[Classical L-function contour]\label{rem:mellin}
Let $f\in S_k$. As $\tau_0\to i\infty$, the $S$-segment maps onto the standard 
L-function contour, $[0,i\infty]$. Simultaneously, as $f(it) \sim e^{-2\pi t}$ as 
$t \to \infty$, the $T$-segment vanishes exponentially. Thus, in this limit, $L^*(f,s)$ 
exactly maps onto the classical L-function integral~\cite{01-Apostol}:
\begin{equation}
\lim_{\tau_0\to i\infty}L^*(f,s)=\int_0^{\infty}\! t^{s-1}f(it)\,dt.
\end{equation}
\end{remark}

\section{$L$-functions for weak modular forms}\label{sec:Lfun}% TODO placeholder title --- no person-names yet

\begin{proposition}[Hurwitz-Gamma]\label{prop:HGG}
For $n\ne0$, $\nu\in\CC$, and $z\in\HH$,
\begin{equation}
\int_{z}^{z+1}\! e^{2\pi inu}\,\zeta(\nu,u)\,du=e^{i \pi (1-\nu)/2} (2\pi n)^{\nu-1}\,\Gamma\big(1-\nu,-2\pi i n z\big).
\end{equation}
At $z=i$ the argument $-2\pi i n z$ collapses to $2\pi n$, giving $\Gamma(1-\nu,2\pi n)$.
\end{proposition}

\begin{proof}
Write $\zeta(\nu,u)=\sum_{m\ge0}(m+u)^{-\nu}$ and shift $u\mapsto u+m$ in the $m$-th term; as
$e^{2\pi in(u+m)}=e^{2\pi inu}$ for integer $m$, the unit segments stack into one horizontal ray,
\begin{equation}
\int_z^{z+1}\! e^{2\pi inu}\zeta(\nu,u)\,du
=\sum_{m\ge0}\int_{z+m}^{z+1+m}\! e^{2\pi inv}v^{-\nu}\,dv
=\int_z^{z+\infty}\! e^{2\pi inv}v^{-\nu}\,dv.
\end{equation}
Substitute $w=-2\pi i n v$, so that $v=iw/(2\pi n)$, $dv=i\,dw/(2\pi n)$, and $e^{2\pi inv}=e^{-w}$:
the lower endpoint becomes $w=-2\pi i n z$, and for $n>0$ the ray rotates to $w\to+\infty$ along the
positive reals --- legal, since $e^{-w}w^{-\nu}$ decays in the right half-plane, so the closing arc
vanishes. With $v^{-\nu}=i^{-\nu}(2\pi n)^{\nu}w^{-\nu}$,
\begin{equation}
\int_z^{z+\infty}\! e^{2\pi inv}v^{-\nu}\,dv
=i^{1-\nu}(2\pi n)^{\nu-1}\!\int_{-2\pi i n z}^{\infty}\! e^{-w}w^{-\nu}\,dw
=e^{i \pi (1-\nu)/2}(2\pi n)^{\nu-1}\Gamma\big(1-\nu,-2\pi i n z\big).
\end{equation}
For $n<0$ the ray instead grows as $\im v\to\infty$ and the rotation is illegal; but both sides are
analytic in $n$ and agree for $n>0$, so the identity continues to $n<0$ with $\Gamma$ on its principal
branch. At $n=0$ the prefactor $(2\pi n)^{\nu-1}$ is singular --- that mode, the constant Fourier
coefficient, is handled separately.
\end{proof}

\begin{theorem}[$L$-function for weakly holomorphic forms]\label{thm:weakL}
For $f\in M^!_k$, we may set $\tau_0=i$, and prove
\begin{equation}\label{eq:weakL}
L^*(f,s)=-c_f(0)\left(\frac1s+\frac{i^k}{k-s}\right)
+\sum_{n\ne0}c_f(n)\left(\frac{\Gamma(s,2\pi n)}{(2\pi n)^s}+i^k\frac{\Gamma(k-s,2\pi n)}{(2\pi n)^{k-s}}\right)~.
\end{equation}
This L-integral satisfies the functional equation $L^*(f,s) = i^k L^*(f,k-s)$.
\end{theorem}

\begin{proof}
By path independence and base-point independence from Proposition~\ref{prop:indep}, 
we may freely set $\tau_0=i$. Here, the $S$-segment degenerates ($-1/i=i$), and 
the L-integral simplifies dramatically:
\begin{equation}
L^*(f,s)=e^{-\pi i s/2}\int_{i-1}^{i} f\,\tk\,d\tau~. 
\end{equation}
We then expand $f$ in terms of its q-series, $\sum_n c_f(n)e^{2\pi in\tau}$,
and substitute $u=\tau+1$, and note that $e^{2\pi i n(u-1)}=e^{2\pi i n u}$. 
Here, each $n\neq0$ mode is handled by  Proposition~\ref{prop:HGG}. 
The first kernel term with $\nu=1-s \leftrightarrow 1-\nu=s$ gives 
\begin{equation}
e^{-\pi i s/2}\,e^{\pi i s/2}(2\pi n)^{-s}\Gamma(s,2\pi n)=\Gamma(s,2\pi n)/(2\pi n)^{s}~,
\end{equation}
and the second kernel term with $\nu=s-k+1 \leftrightarrow 1-\nu=k-s$ gives
\begin{equation}
-\,e^{-\pi i s/2}e^{i\pi(s-1)}\,e^{\pi i (k-s)/2}\,(2\pi n)^{s-k}\,\Gamma(k-s,2\pi n)
=i^{k}\,\frac{\Gamma(k-s,2\pi n)}{(2\pi n)^{k-s}}.
\end{equation}
So each $n\ne0$ mode contributes $\Gamma(s,2\pi n)/(2\pi n)^{s}+i^{k}\Gamma(k-s,2\pi n)/(2\pi n)^{k-s}$. 
Crucially, the $n=0$ mode contributes 
\begin{equation}
c_f(0)\,e^{-\pi i s/2}\!\int_{i-1}^{i}\tk\,d\tau=-c_f(0)\left(\frac1s+\frac{i^k}{k-s}\right)~.
\end{equation}
Combining these reproduces Eq.~\eqref{eq:weakL}. Finally, the functional equation, $L^*(f,s)=i^k L^*(f,k-s)$, follows by direct inspection.
\end{proof}

\subsection{Explicit S- and T-contours for general $\tau_0$}

\begin{lemma}[$T$-segment closed form on $M^!_k$]\label{lem:Tclosed-Mk}
For $f\in M^!_k$, for generic $\tau_0 \in \HH$ away from cusps, denote the T-contour's
contribution to $L^*(f, s)$ as $L^*_T(f, s; \gamma^T_{\tau_0})$. For generic $\tau_0$, we have:
\begin{equation}\label{eq:T-segment-closed-Mk}
\begin{split}
L_T^*(f, s; \gamma^T_{\tau_0})
&= -c_f(0)\!\left(\frac{e^{-\pi i s/2}\tau_0^{\,s}}{s} + \frac{e^{\pi i s/2}\tau_0^{\,k-s}}{k-s}\right)\\
&\quad + \sum_{n \ne 0}\! c_f(n)\!\left(\frac{\Gamma(s, -2\pi i n\tau_0)}{(2\pi n)^s} + i^k\frac{\Gamma(k - s, -2\pi i n\tau_0)}{(2\pi n)^{k-s}}\right).
\end{split}
\end{equation}
At $\tau_0 = i$ the constant mode collapses to $-c_f(0)(1/s + i^k/(k-s))$, recovering Theorem~\ref{thm:weakL}.
\end{lemma}
\begin{proof}
Expand $f = \sum_n c_f(n)\,q^n$ on the compact segment
$[\tau_0 - 1, \tau_0]$, and restrict attention to the sub-sum, $\sum_{n \neq 0}$. 
For each such $n \neq 0$, apply the Proposition~\ref{prop:HGG} in
$\widetilde{k}_T(\tau, s) = \zeta(1 - s, \tau + 1) - e^{i\pi(s-1)}\zeta(s - (k-1), \tau + 1)$
at $\nu = 1 - s$ and $\nu = s - (k - 1)$. The phases for the second term work-out 
identically as in the proof of Theorem~\ref{thm:weakL}.  

For $n = 0$, the constant $c_f(0)$ multiplies $e^{-\pi i s/2}\!\int_{\tau_0 - 1}^{\tau_0}\!\widetilde k_T(\tau, s)\,d\tau$;
the Hurwitz antiderivative $\int_{\tau_0 - 1}^{\tau_0}\zeta(\sigma, \tau + 1)\,d\tau = \tau_0^{1 - \sigma}/(\sigma - 1)$
--- from $\partial_\tau\zeta(\sigma - 1, \tau + 1) = -(\sigma - 1)\zeta(\sigma, \tau + 1)$ and the shift
$\zeta(\sigma - 1, \tau_0 + 1) - \zeta(\sigma - 1, \tau_0) = -\tau_0^{1 - \sigma}$ --- applied to the two
halves of $\widetilde k_T$ at $\sigma = 1 - s$ and $\sigma = s - k + 1$ yields
$-e^{-\pi i s/2}\tau_0^{\,s}/s$ and $-e^{\pi i s/2}\tau_0^{\,k-s}/(k-s)$ respectively.
\end{proof}


\begin{corollary}[$T$-segment cusp asymptotics]\label{lem:cusp-decay}
Let $f \in S^!_k$, where $f = \sum_{n \geq -N} c_f(n)q^n$ for some $-N \in \ZZ$. 
The T-contour's contribution to $L^*(f,s)$ is given by 
Lemma~\ref{lem:Tclosed-Mk}. As $\tau_0 = it$ goes to the cusp $i\infty$,
the $T$-contour scales exactly as $f(it)$ scales:
$L_T^*(f, s; \gamma^T_{\tau_0}) \sim f(it) \sim e^{2\pi N t}$.
\end{corollary}
\begin{proof}
For the asymptotics, $\Gamma(\sigma, z) \sim z^{\sigma - 1}e^{-z}$
as $|z| \to \infty$ for $|\arg z| < \tfrac{3\pi}{2}$ \cite[\S8.11]{DLMF}
gives $|c_f(n)|$-weighted go as $e^{-2\pi n\,\im(\tau_0)}$.  The smallest 
mode $n$ in the q-series, i.e. $n = -N$, dominates the scaling of $f(it)$ and of 
$L^*_T(f, s; \gamma^T_{i t})$. Combining these results yields the claimed asymptotics.
\end{proof}

Now, exploiting $\tau_0$-independence of $L^*(f, s)$, its explicit value as given in 
Theorem~\ref{thm:weakL}, and the explicit T-contribution $L^*_T(f, s; \gamma^T_{\tau_0})$ trivially 
implies clean, closed-form expressions for the S-contour's contribution to $L^*(f,s)$: 

\begin{corollary}[$S$-segment closed form]\label{cor:Sclosed}
For $f\in M^!_k$ and $\tau_0\in\HH$, the $S$-contribution
$L_S^*(f,s;\gamma^S_{\tau_0}) := e^{-\pi i s/2}\!\int_{-1/\tau_0}^{\tau_0}\! f(\tau)\,\tau^{s-1}\,d\tau$
is the full $L$-value minus the $T$-contribution at the same base point, hence an explicit difference
of incomplete gammas between the base points $i$ and $\tau_0$:
\begin{align}
L_S^*(f,s;\gamma^S_{\tau_0})
&= L^*(f,s) - L_T^*(f,s;\gamma^T_{\tau_0}) \label{eq:S-segment-closed}\\
&= c_f(0)\left(\frac{e^{-\pi i s/2}\tau_0^{\,s}-1}{s}
+ \frac{e^{\pi i s/2}\tau_0^{\,k-s}-i^k}{k-s}\right) \label{eq:S-segment-explicit}\\
&+ \sum_{n\ne0} c_f(n)\!\left(\frac{\Gamma(s,2\pi n)-\Gamma(s,-2\pi i n\tau_0)}{(2\pi n)^s}
+ i^k\,\frac{\Gamma(k-s,2\pi n)-\Gamma(k-s,-2\pi i n\tau_0)}{(2\pi n)^{k-s}}\right)~. \nonumber
\end{align}
\end{corollary}
\begin{proof}
Line \eqref{eq:S-segment-closed} is the definition $L^*(f,s;\gamma^S_{\tau_0},\gamma^T_{\tau_0})
= L_S^*(\gamma^S_{\tau_0}) + L_T^*(\gamma^T_{\tau_0})$ together with the $\tau_0$-independence of $L^*$
(Proposition~\ref{prop:indep}); evaluating at $\tau_0=i$, where $L_S^*(\gamma^S_i)=0$, identifies
$L^*(f,s)=L_T^*(\gamma^T_i)$. Line \eqref{eq:S-segment-explicit} then inserts the closed forms of
Theorem~\ref{thm:weakL} for $L^*(f,s)$ and Lemma~\ref{lem:Tclosed-Mk} for $L_T^*(\gamma^T_{\tau_0})$:
the $n\ne0$ modes subtract base point by base point into the incomplete-gamma differences, and the two
constant modes subtract to the $c_f(0)$-terms in Eq.~\eqref{eq:S-segment-explicit}.
%\begin{equation}
%-c_f(0)\Big(\frac1s+\frac{i^k}{k-s}\Big)
%+ c_f(0)\Big(\frac{e^{-\pi i s/2}\tau_0^{\,s}}{s}+\frac{e^{\pi i s/2}\tau_0^{\,k-s}}{k-s}\Big)
%= c_f(0)\Big(\frac{e^{-\pi i s/2}\tau_0^{\,s}-1}{s}+\frac{e^{\pi i s/2}\tau_0^{\,k-s}-i^k}{k-s}\Big).
%\end{equation}
\end{proof}


\subsection{Classical Mellin recovery and canonical uniqueness.}\label{sec:Connection-w-previous-weak-MF-l-functions}
As discussed in Remark~\ref{rem:mellin}, for $f \in S_k$, taking the cusp-limit
$\tau_0 \to i\infty$ in Theorem~\ref{thm:weakL} collapses the $S$-segment 
onto the classical Mellin contour $[0, i\infty]$ while the $T$-segment vanishes 
exponentially, since $f(it) = O(e^{-2\pi t})$ on the segment as $\tau = it \to i\infty$. I.e., here, one recovers the classical L-integral of $f$ without obstruction:
\begin{equation}\label{eq:Mellin-recovery}
\lim_{\tau_0 \to i\infty}\!
L^*(f, s; \gamma^S_{\tau_0}, \gamma^T_{\tau_0})
\;=\; \int_0^\infty\! t^{s - 1}\,f(it)\,dt~.
\end{equation}  

Of course, path- and base-point independence of $L^*(f, s) = L^*(f, s; \gamma^S_{\tau_0}, \gamma^T_{\tau_0})$ 
is guaranteed by Theorem~\ref{thm:weakL} for all $f \in M_k^!$. However, these 
L-integrals are the sums of the S-contour and T-contour. Finiteness is ensured only 
of the sum of these two parts. Both contours may, and indeed do, grow without bound 
as $\tau_0 \to i \infty$ when $f \notin M_k$, as is highlighted in Corollary~\ref{lem:cusp-decay}:

As $\tau_0\to i\infty$ the polar mode $q^{-N}$ of a properly weakly holomorphic modular form 
$f \in M_k^! $ makes the $T$-contribution grow like $e^{2\pi N\,\im\tau_0}$ and the $S$-contribution
diverge oppositely. Yet their sum, the $\tau_0$-independent value $L^*(f,s)$, remains finite and constant
over all finite $\tau_0 \in \HH$. Discarding $\gamma^T_{\tau_0}$ therefore does violence to the result: 
When the $S$-contour is isolated from its T-partner, and $\tau_0 \to i \infty$, it diverges. This is the precise
reason that~\cite{BFK2011} needed to regularize their divergent L-integral representation for properly weak cusp forms.

\subsection{Period polynomials, periods and quasi-periods}

Because $L^*(f,s)$ is canonical for every $f\in M^!_k$, the period polynomial needs no separate
construction --- it is the generating polynomial of the critical values.

\begin{definition}[Period polynomial]\label{def:rf}
For $f\in S^!_k$ and $n=k-2$, the \emph{period polynomial} is
\begin{equation}
r_f(X,Y):=(2\pi i)^{n+1}\sum_{\ell=0}^{n}(-1)^\ell\binom{n}{\ell}\,i^{\ell+1}\,L^*(f,\ell+1)\,X^{n-\ell}Y^\ell.
\end{equation}
\end{definition}

\begin{proposition}[Period polynomial as a contour integral]\label{prop:periodint}
With $(X-\tau Y)^n$ in place of the $L$-weight $\tau^{s-1}$, and
$\mathbf K_T(\tau;X,Y):=\sum_{\ell=0}^{n}\binom{n}{\ell}(-1)^\ell\,\tk(\tau,\ell+1)\,X^{n-\ell}Y^\ell$
in place of $\tk$, the same two segments assemble $r_f$ directly:
\begin{equation}
r_f=(2\pi i)^{n+1}\Big[\int_{\gamma^S} f(\tau)\,(X-\tau Y)^n\,d\tau
+\int_{\gamma^T} f(\tau)\,\mathbf K_T(\tau;X,Y)\,d\tau\Big].
\end{equation}
\end{proposition}
\begin{proof}
By Definition~\ref{def:Lint}, the $X^{n-\ell}Y^\ell$ coefficient of the bracket is
\begin{equation}
r_f|_{X^{n-\ell}Y^\ell} = (-1)^\ell\binom{n}{\ell}\big[\int_{\gamma^S} f\,\tau^\ell+\int_{\gamma^T} f\,\tk(\tau,\ell+1)\big]
=(-1)^\ell\binom{n}{\ell}\,i^{\ell+1}L^*(f,\ell+1)~.
\end{equation}
This matches definition~\ref{def:rf} term by term, and proves the claim.
\end{proof}

\paragraph{Quasi-periods and their determinant:} 
As an aside, the Hurwitz T-kernel in this note was originally uncovered in order to reproduce 
the manifestly finite computations of the quasi-periods for $\Dhat = 1/q + O(q^2) \in S_{12}^!$
originally given in~\cite{Brown2017}. In this vein, it is interesting to show collect quasi-periods 
computable from our L-integrals/period-polynomials for $f \in S_k^!$ (when $\dim(S_k) = 1$).

To this effect, we recall the notation of~\cite{DukeJenkins}. Decomposing the modular weight 
$k = 12 \lambda + k'$ where $k' \in \{ 0, 4, 6, 8, 10, 14\}$, we see $\lambda = \dim(S_k)$ when 
$\lambda \geq 0$. For any $\lambda$, we may denote $f_{k, m} := 1/q^m + O(q^{\lambda+1}) \in M_k^!$. 
When $\dim(S_k) = 1$, we have $f_{k,m} = 1/q^m + O(q^2)$. Here, define 
$\Delta_k := f_{k,-1} = q + O(q^2) \in S_k$ and $\Dhat_k := f_{k,1} = 1/q + O(q^2) \in S_k^!$.
Further, define
\begin{align}
{\rm pd}_k := \left( \begin{smallmatrix} 
\omega^+(\Delta_k) & \omega^-(\Delta_k) \\
\eta^+(\Dhat_k) & \eta^-(\Dhat_k)
\end{smallmatrix}\right)
\end{align}
That $\det({\rm pd}_k)$ satisfies a Legendre-type relation, $\det({\rm pd}_k)\in\QQ^\times\cdot(2\pi i)^{k-1}$,
is the de Rham theory of Brown--Hain~\cite[Prop.~5.6]{BrownHain2018}. Brown~\cite[\S8.3]{Brown2017}
evaluates the constant at $k=12$, finding $\det({\rm pd}_k)=-(k-2)!\,(2\pi i)^{k-1}$ --- his period
matrix orders the quasi-periods first, giving $+(k-2)!$, where ours orders the periods first --- and
remarks that he could locate no reference for the value. We reproduce $k=12$ and give the analogous
rational multiples for $k \in \{16, 18, 20, 22, 26\}$ in Table~\ref{tab:periods}, below:

\begin{table}[H]\centering\small
\begin{tabular}{r|rrrr|c}
$k$ & $\omega^+(f_k)$ & $\omega^-(f_k)/i$ & $\eta^+(\Dhat_k)$ & $\eta^-(\Dhat_k)/i$ & $\frac{-i^k\det({\rm pd}_k)}{(2\pi)^{k-1}(k-2)!}$\\\hline
12 & $-6.89168\mathrm{e}7$  & $-5.58502\mathrm{e}6$  & $1.27202\mathrm{e}{11}$ & $1.02767\mathrm{e}{10}$  & $1$\\
16 & $-8.81481\mathrm{e}{11}$ & $-1.54063\mathrm{e}{10}$ & $4.14120\mathrm{e}{13}$ & $6.83541\mathrm{e}{11}$ & $13/30$\\
18 & $-5.07833\mathrm{e}{13}$ & $-5.43305\mathrm{e}{12}$ & $-2.19200\mathrm{e}{14}$ & $-3.36339\mathrm{e}{13}$ & $2/3$\\
20 & $-4.22882\mathrm{e}{16}$ & $-3.73170\mathrm{e}{14}$ & $6.38640\mathrm{e}{16}$ & $-1.89823\mathrm{e}{14}$ & $17/5$\\
22 & $-1.49974\mathrm{e}{18}$ & $-4.24489\mathrm{e}{16}$ & $1.43795\mathrm{e}{18}$ & $-4.41050\mathrm{e}{16}$ & $19/21$\\
26 & $-1.18673\mathrm{e}{22}$ & $-3.60961\mathrm{e}{21}$ & $1.18653\mathrm{e}{22}$ & $-3.61016\mathrm{e}{21}$ & $23/15$\\
\end{tabular}
\caption{
Periods $\omega^\pm(\Delta_k)$ and 
quasi-periods $\eta^\pm(\Dhat_k)$, and the rational number
$(-1)^{1+k/2}\det/\big((2\pi i)^{k-1}(k-2)!\big)$ with
$\det({\rm pd}_k) =\omega^+\eta^--\omega^-\eta^+$.
At $k=12$ the value matches Brown~\cite{Brown2017}; the entries for $k\ge16$ appear to be new.
These ratios are not intrinsic to $f$: they depend on how the rational period-polynomial basis
$P^\pm$ is normalized (we use a primitive-integer normalization; Brown's normalization gives
$-(k-2)!$ at $k=12$). We have not found a closed form for them.
Full 50-digit values in~\cite{VerifyScripts}.
}\label{tab:periods}
\end{table}


\section{Meromorphic modular forms}\label{sec:mero}% TODO placeholder title

We find it convenient to define the following function,
\begin{align}\label{eq:gc}
G_{s, k}(n) := \frac{\Gamma(s, 2\pi n)}{(2\pi n)^s} + i^k \frac{\Gamma(k-s, 2 \pi n)}{(2\pi n)^{k-s}}.
\end{align}


\begin{proposition}[Basepoint and path independence within a homotopy class]\label{prop:indepHomotopy}
For $f\in F_k$, the value $L^*(f,s;\gamma^S_{\tau_0},\gamma^T_{\tau_0})$ is independent of $\tau_0$
and of the path so long as both vary within a single homotopy class of admissible contours in
$\HH\setminus\Gamma\!\cdot\!f^{-1}(\infty)$. For $f\in M^!_k$ there is a single class, recovering
Proposition~\ref{prop:indep}; across different classes the value changes by the residues the
contour sweeps, as quantified in Corollary~\ref{cor:wall}.
\end{proposition}

\begin{proof}
The $\tau_0$-derivative of Proposition~\ref{prop:indep} uses only the $T$-periodicity
$f(\tau_0-1)=f(\tau_0)$ and the modularity $f(-1/\tau_0)=\tau_0^k f(\tau_0)$, both valid for every
$f\in F_k$, and is evaluated at the moving endpoints alone. While $\tau_0$ stays in one chamber of
$\HH\setminus f^{-1}(\infty)$ --- where the segments meet no pole --- the boundary terms cancel
exactly as there, so $\partial_{\tau_0}L^*=0$; the interior poles lie off the contour and never
enter. Within a fixed class the integrands $f\,\tau^{s-1}$ and $f\,\tk$ are holomorphic along the
contours, so any two are joined by a deformation avoiding $\Gamma\!\cdot\!f^{-1}(\infty)$, and
Cauchy's theorem equates their integrals.
\end{proof}


% §5.14a — [Def] polylog projection (def:proj)
\begin{definition}[Polylog projection]\label{def:proj}
Let $f\in F_k$ have interior poles at the points $\tau_p$ of a set $\mathcal P_f$ of
$\Gamma$-inequivalent representatives, and write $\Li_s(z)=\sum_{n>0}z^n n^{-s}$. Since
$\Li_{1-m}\big(e^{2\pi i(\tau-\tau_p)}\big)$ has a pole of order $m$ at $\tau_p$, the principal part of
$f$ there is matched by a unique \emph{polylog projection}
\begin{equation}
\widehat P_{\tau_p}f(\tau):=\sum_{m>0}r^*_{f,\tau_p}(m)\,\Li_{1-m}\big(e^{2\pi i(\tau-\tau_p)}\big),
\end{equation}
its coefficients $r^*_{f,\tau_p}(m)$ being the principal-part data of $f$ at $\tau_p$. The
\emph{projection} $\widehat P:=\sum_{\im\tau_p\ge1}\widehat P_{\tau_p}$ removes the principal
part of every pole on or above the line $\im\tau=1$.
\end{definition}

% §5.14b — [Def] regularized cusp coefficients (def:cRf)
\begin{definition}[Regularized cusp coefficients]\label{def:cRf}
With $\widehat P$ as in Definition~\ref{def:proj},
\begin{equation}
c_{Rf}(n):=\big[(I-\widehat P)f\big]_{q^n}.
\end{equation}
\end{definition}

\begin{definition}[The space $F(\infty)$]\label{def:Finfty}
$F(\infty)$ is the space of meromorphic $g$ on $\HH$ with $g(\tau)=g(\tau+1)$, finitely many
poles in any compact subset of the strip $\{{\rm Im}\,\tau>0,\ |{\rm Re}\,\tau|\le\tfrac12\}$,
and $g(\tau)=O(e^{C\,{\rm Im}\,\tau})$ as ${\rm Im}\,\tau\to\infty$ for some $C$. Each such
$g$ has a Fourier expansion $g=\sum_n c_g(n)q^n$.
\end{definition}

Then $F_k\subset F(\infty)$, and $(I-\widehat P)f\in F(\infty)$ --- the polylog tails are
$T$-periodic but not modular, so the projection leaves $F_k$.

\begin{lemma}[$q$-series--integral commutation]\label{lem:commute}
Let $g\in F(\infty)$ have no pole on or above the line ${\rm Im}\,\tau=1$ (e.g.\
$g=(I-\widehat P)f$). Then $g=\sum_n c_g(n)q^n$ converges along the $T$-segment at $\tau_0=i$
and integrates term by term:
\begin{equation}
\int_{\gamma^T_i} g(\tau)\,\tk(\tau,s)\,d\tau
=\sum_n c_g(n)\int_{i-1}^{i} e^{2\pi in\tau}\,\tk(\tau,s)\,d\tau.
\end{equation}
A pole at ${\rm Im}\,\tau_p\ge1$ defeats this --- there the coefficients grow at least like
$n^{m-1}e^{2\pi n\,{\rm Im}\,\tau_p}$~\cite[Lem.~2.6]{1806}.
\end{lemma}

\begin{proof}
With no pole on or above ${\rm Im}\,\tau=1$, the finitely many poles of $g$ in the strip have
heights $<1$; let $h<1$ be the largest. The Fourier series $g=\sum_n c_g(n)q^n$ converges
absolutely and uniformly on compact subsets of $\{{\rm Im}\,\tau>h\}$, in particular on the
$T$-segment $[i-1,i]$; with $\tk(\tau,s)$ bounded there, the sum passes through the integral,
which is the displayed identity. Conversely a pole at height ${\rm Im}\,\tau_p\ge1$ forces
$|c_g(n)|\gtrsim n^{m-1}e^{2\pi n\,{\rm Im}\,\tau_p}$~\cite[Lem.~2.6]{1806}, so
$|c_g(n)q^n|\gtrsim n^{m-1}e^{2\pi n({\rm Im}\,\tau_p-1)}\not\to0$ on $[i-1,i]$ and the series
diverges.
\end{proof}

\begin{definition}[Polylog integral]\label{def:polylog}
For $N\ge0$ and $z\in\HH$, the \emph{polylog integral} is the $L$-integral of a single tail
$\Li_{-N}(e^{2\pi i(\tau-z)})$ along the $T$-segment at $\tau_0=i$ (where $\gamma^S$ degenerates, since
$-1/i=i$):
\begin{equation}
I_N(s,z):=e^{-\pi i s/2}\int_{i-1}^{i}\Li_{-N}\big(e^{2\pi i(\tau-z)}\big)\,\tk(\tau,s)\,d\tau.
\end{equation}
\end{definition}

\begin{lemma}[Specialized polylog integral]\label{lem:polylogcf}
Recall $G_{s, k}(n)$ from Eq.~\eqref{eq:gc}. For the projected poles
${\rm Im}\,z>1$, where $|e^{2\pi i(\tau-z)}|\ge1$ on the contour, polylog inversion
$\Li_{-N}(w)=(-1)^{N+1}\Li_{-N}(1/w)$ resums the integral into the convergent series
\begin{equation}
I_N(s,z)=\delta_{N,0}\Big(\tfrac1s+\tfrac{i^k}{k-s}\Big)
+(-1)^{N+1}\sum_{n\ge1} n^N e^{2\pi inz}\,G_{s,k}(-n).
\end{equation}
This sum converges exponentially quickly, with base $e^{-2\pi n (y -1)}$, where $y = {\rm Im}(z)$.
\end{lemma}

\begin{proof}
For ${\rm Im}\,z>1$, $|e^{2\pi i(\tau-z)}|\ge1$ on $[i-1,i]$, so the series of $\Li_{-N}(e^{2\pi i(\tau-z)})$
diverges; the inversion rewrites it through $\Li_{-N}(e^{2\pi i(z-\tau)})=\sum_{n\ge1}n^N e^{2\pi in(z-\tau)}$,
convergent since $|e^{2\pi i(z-\tau)}|\le1$. Integrating term by term against $\tk$, each
$e^{-\pi i s/2}\!\int_{i-1}^i e^{-2\pi in\tau}\tk\,d\tau$ is the building block $G_{s,k}(n)$ of
Theorem~\ref{thm:weakL} at mode $-n$, giving the displayed sum; for $N=0$ the leftover $-1$ in
$\Li_0(w)=-1-\Li_0(1/w)$ adds $-e^{-\pi i s/2}\!\int_{i-1}^i\tk=\tfrac1s+\tfrac{i^k}{k-s}$. 

Exponential convergence $\sim e^{-2\pi n(y-1)}$ for $y = {\rm Im}(z) > 1$ follows from combining the asymptotic 
scaling $\Gamma(s,-2\pi n) \sim (2\pi n)^{s-1} e^{2\pi n}$ and the exact result $|e^{2\pi i n z}| = e^{-2\pi n y}$.
\end{proof}

% [Former lem:polylogEdge / lem:polylogCorner (edge z=i+x and corner z=i) removed in the §4
%  reorg: a pole on the line Im(tau)=1 sits ON the contour, so it is an ordinary on-contour
%  pole --- principal value, with the residue jump governed by Corollary~\ref{cor:wall} ---
%  not a projected pole needing its own closed form.  L* for such a form is basepoint-
%  independent (numerically verified, incl. the double pole of E_4 Delta/E_6^2 at tau=i).]

\begin{theorem}[$L$-function of a meromorphic modular form]\label{thm:mero}
For $f\in F_k$ with no pole on the edge ${\rm Im}\,\tau=1$, at $\tau_0=i$,
\begin{equation}
L^*(f,s)=-c_{Rf}(0)\Big(\tfrac1s+\tfrac{i^k}{k-s}\Big)
+\sum_{n\ne0}c_{Rf}(n)\,G_{s, k}(n)
+\sum_{{\rm Im}\,\tau_p>1}\sum_{m>0}r^*_{f,\tau_p}(m)\,I_{m-1}(s,\tau_p),
\end{equation}
where, again $G_{s, k}(n)$ is the sum of incomplete gammas defined in Eq.~\eqref{eq:gc};
the regularized cusp sum (Definition~\ref{def:cRf}, Theorem~\ref{thm:weakL}) plus one polylog
integral (Lemma~\ref{lem:polylogcf}) per interior pole above the edge. A pole on the edge
${\rm Im}\,\tau=1$ --- the elliptic point $\tau=i$ included --- instead sits on the contour: there
$L^*$ is the principal value, and sweeping the contour across the pole changes it by the residue
of Corollary~\ref{cor:wall}.
\end{theorem}

\begin{proof}
At $\tau_0=i$ the $S$-segment degenerates, so $L^*(f,s)=e^{-\pi i s/2}\int_{i-1}^i f\,\tk\,d\tau$. Split
$f=(I-\widehat P)f+\widehat Pf$. The regularized part has no pole on or above ${\rm Im}\,\tau=1$
(Definition~\ref{def:proj}), so by Lemma~\ref{lem:commute} its Fourier series integrates term by
term; mode by mode this is Theorem~\ref{thm:weakL}, giving
$-c_{Rf}(0)(\tfrac1s+\tfrac{i^k}{k-s})+\sum_{n\ne0}c_{Rf}(n)G_{s,k}(n)$. The projected part
$\widehat Pf=\sum_{\tau_p}\sum_{m>0}r^*_{f,\tau_p}(m)\Li_{1-m}(e^{2\pi i(\tau-\tau_p)})$ integrates to
$\sum_{{\rm Im}\,\tau_p>1}\sum_{m>0}r^*_{f,\tau_p}(m)I_{m-1}(s,\tau_p)$ by
Lemma~\ref{lem:polylogcf}. Their sum is the stated formula.
\end{proof}

%DAM: I think this lemma is not needed. I think we can do a generalization of~\ref{lem:polylogcf} 
%where we simply analytically continue the closed-form result to the special cases of $\tau_0 = i$ 
%or $\tau_0) = i + x$ for $x \in \RR$. 
%\begin{lemma}[$L$-function, a pole on the edge]\label{lem:onedge}
%Routing the $T$-path just below versus just above a pole $\tau_p$ of $f\in F_k$ at height
%${\rm Im}\,\tau_p=1$ gives finite values differing by the residue,
%\begin{equation}
%L^*_{\rm below}(f,s)-L^*_{\rm above}(f,s)
%=2\pi i\,e^{-\pi i s/2}\Res_{\tau=\tau_p}\!\big[f(\tau)\,\tk(\tau,s)\big].
%\end{equation}
%We define the edge value as the principal value $L^*=\tfrac12(L^*_{\rm above}+L^*_{\rm below})
%=L^*_{\rm below}-\pi i\,e^{-\pi i s/2}\Res_{\tau=\tau_p}[f\,\tk]$. At the elliptic point $\tau_p=i$ the
%pole lies on the $S$-segment instead; its principal value, reached as $\tau_0\to i$, is finite
%and basepoint-independent.
%\end{lemma}
%
%\begin{proof}
%The two routings of $\gamma^T$ differ by a contour encircling $\tau_p$ once, so the residue
%theorem gives $L^*_{\rm below}-L^*_{\rm above}=2\pi i\,e^{-\pi i s/2}\Res_{\tau=\tau_p}[f\,\tk]$, and the
%principal value is their mean. At $\tau_p=i$ the pole sits on the degenerate $S$-segment;
%displacing to $\tau_0=i(1+\delta)$ keeps it on $\gamma^S$ for all small $\delta>0$, so no wall is
%crossed and $L^*$ is finite and $\delta$-independent (verified numerically~\cite{VerifyScripts};
%a closed proof is the $S$-segment principal value of~\cite{1806}).
%\end{proof}

% DAM: rewritten per your margin note --- the S-segment carries its own windings
% (rem:Snonvanish). Please vet the X_S,X_T framing and the reference-class convention.
\begin{corollary}[Wall-crossing and the winding data $X_S,X_T$]\label{cor:wall}
Fix $f\in F_k$ and record, for an admissible pair $(\gamma^S_{\tau_0},\gamma^T_{\tau_0})$, its
\emph{winding data} --- functions $X_S,X_T:\Gamma\!\cdot\!f^{-1}(\infty)\to\ZZ$ giving the net number of
turns of $\gamma^S$, resp.\ $\gamma^T$, about each pole relative to a fixed reference pair (value
$L^*_0(f,s)$ there). By Proposition~\ref{prop:indepHomotopy} the integral depends only on $(X_S,X_T)$:
\begin{equation}\label{eq:wall-crossing}
L^*(f,s;X_S,X_T) = L^*_0(f,s) + 2\pi i\,e^{-\pi i s/2}\!\left[
\sum_{\tau_p} X_S(\tau_p)\,\Res_{\tau=\tau_p}\!\big(f\,\tau^{s-1}\big)
+ \sum_{\tau_p} X_T(\tau_p)\,\Res_{\tau=\tau_p}\!\big(f\,\tk\big)\right].
\end{equation}
A single wall-crossing is the rank-one case: sweeping one segment once across one pole $\tau_p$ shifts
$L^*$ by $2\pi i\,e^{-\pi i s/2}\Res_{\tau_p}(f\,\tau^{s-1})$ for $\gamma^S$, or
$2\pi i\,e^{-\pi i s/2}\Res_{\tau_p}(f\,\tk)$ for $\gamma^T$. Thus $L^*$ is fixed by the contour only up
to the residues its two segments wind around --- weighted by $\tau^{s-1}$ on $\gamma^S$ and by $\tk$ on
$\gamma^T$.
\end{corollary}
\begin{proof}
Within-class constancy is Proposition~\ref{prop:indepHomotopy}. Two pairs whose winding data differ by
one unit at a single $\tau_p$ on one segment bound a contour encircling $\tau_p$ once; the residue
theorem adds $2\pi i$ times the residue of that segment's integrand --- $f\,\tau^{s-1}$ for $\gamma^S$,
$f\,\tk$ for $\gamma^T$ --- and the prefactor $e^{-\pi i s/2}$ of Definition~\ref{def:Lint} multiplies
through. Summing over all crossings gives \eqref{eq:wall-crossing}.
\end{proof}

\begin{remark}[The $S$-segment does not drop out]\label{rem:Snonvanish}
For $f\in M^!_k$ the interior is pole-free, so $X_S\equiv X_T\equiv0$ and $L^*$ is path-independent
(Proposition~\ref{prop:indep}). For $f\in F_k$ both windings are live, and the $S$-segment is no
exception. Its endpoints $-1/\tau_0$ and $\tau_0$ are $S$-equivalent, so $\gamma^S$ can be closed up
through the group action into a loop --- but a loop enclosing interior poles does \emph{not} integrate
to zero in the meromorphic case; it equals $2\pi i\sum(\text{residues})$. The identification of the
$S$-endpoints therefore does not make $\gamma^S$ negligible: it carries genuine residue data $X_S$,
weighted by $\tau^{s-1}$, on exactly the same footing as the $T$-segment's $X_T$ weighted by $\tk$.
\end{remark}

\begin{thebibliography}{99}

\bibitem{BFK2011}
K.\ Bringmann, K.-H.\ Fricke, and Z.\ Kent,
\emph{Special $L$-values and periods of weakly holomorphic modular forms},
Ramanujan J.\ \textbf{24} (2011), 119--139.

\bibitem{Brown2017}
F.\ Brown, \emph{A class of non-holomorphic modular forms III},
arXiv:1710.07912 (2017).

\bibitem{1806}
D.\ A.\ McGady,
\emph{$L$-functions for meromorphic modular forms and sum rules in
conformal field theory},
arXiv:1806.09874 (2018).

\bibitem{01-Apostol}
T. M. Apostol,
``Modular functions and Dirichlet series in number theory,''
Graduate Texts in Mathematics, London, Springer (1997).


\bibitem{DukeJenkins}
W.\ Duke and P.\ Jenkins,
\emph{On the zeros and coefficients of certain weakly holomorphic modular forms},
Pure Appl.\ Math.\ Q.\ \textbf{4} (2008), no.\ 4, 1327--1340.

% Rational period polynomials: classical background/normalization for the
% period-polynomial and quasi-period material.
\bibitem{KZ1984}
W.\ Kohnen and D.\ Zagier,
\emph{Modular forms with rational periods},
in \emph{Modular Forms} (R.~A.\ Rankin, ed.), Ellis Horwood, 1984, pp.\ 197--249.



\bibitem{BrownHain2018}
F.\ Brown and R.\ Hain,
\emph{Algebraic de Rham theory for weakly holomorphic modular forms of level one},
Algebra Number Theory \textbf{12} (2018), no.\ 3, 723--750; arXiv:1708.00652.

\bibitem{VerifyScripts}
\emph{Numerical verification scripts} (companion code), 2026.

\bibitem{DLMF}
NIST Digital Library of Mathematical Functions,
F.~W.~J.\ Olver \emph{et al.}, eds.,
\url{https://dlmf.nist.gov/}, Release 1.2.x.
Chapters used: \S8 (incomplete gamma) and \S25.11 (Hurwitz zeta).

% =====================================================================
% Adjacent work --- not yet \cite'd in the body; intended for the
% (forthcoming) introduction / motivation.  Each squib = intended role.
% =====================================================================

% Eichler--Shimura / automorphic cohomology: the cohomological backdrop for
% reading the two-segment integral as a finite-contour cocycle.
\bibitem{BCD2018}
R.\ Bruggeman, Y.\ Choie, and N.\ Diamantis,
\emph{Holomorphic automorphic forms and cohomology},
Mem.\ Amer.\ Math.\ Soc.\ \textbf{253} (2018), no.\ 1212; arXiv:1404.6718.

% Eichler--Shimura for mock/weak forms: motivates L-functions of weakly
% holomorphic cusp forms (the BGKO ``extra relation'').
\bibitem{BGKO}
K. Bringmann, P. Guerzhoy, Z. Kent and K. Ono,
\emph{Eichler-Shimura Theory for Mock Modular Forms},
Mathematische Annalen, March 2013 355(3).

% Period polynomials and the Haberland pairing: the normalization behind the
% quasi-period determinant (Table~\ref{tab:periods}).
\bibitem{PasolPopa2013}
V.\ Pa\c{s}ol and A.\ A.\ Popa,
\emph{Modular forms and period polynomials},
Proc.\ London Math.\ Soc.\ (3) \textbf{107} (2013), 713--743; arXiv:1202.5802.

% Closest adjacent construction: L-functions of meromorphic (quasi-)modular
% forms --- the natural comparison point for our meromorphic L-function.
\bibitem{WangZhang2022}
W.\ Wang and H.\ Zhang,
\emph{Meromorphic quasi-modular forms and their $L$-functions},
arXiv:2202.09542 (2022).

% Periods of meromorphic modular forms (motivic / single-valued): adjacent
% treatment of the meromorphic case.
\bibitem{BrownFonseca2025}
F.\ Brown and T.\ J.\ Fonseca,
\emph{Single-valued periods of meromorphic modular forms and a motivic
interpretation of the Gross--Zagier conjecture},
arXiv:2508.04844 (2025).

% Regularized inner products: the non-holomorphic regularization our finite
% contour avoids (intro contrast with BFK-type constructions).
\bibitem{BDE2016}
K.\ Bringmann, N.\ Diamantis, and S.\ Ehlen,
\emph{Regularized inner products and errors of modularity},
Int.\ Math.\ Res.\ Not.\ IMRN (2016); arXiv:1603.03056.

% Period functions of real-analytic forms: the Brown non-holomorphic-forms
% context for periods and quasi-periods.
\bibitem{DiamantisDrewitt2020}
N.\ Diamantis and J.\ Drewitt,
\emph{Period functions associated to real-analytic modular forms},
Res.\ Math.\ Sci.\ \textbf{7} (2020), no.\ 21.


%\bibitem{companion}
%  D.\ A.\ McGady,
%  \emph{Quasi-periods, periods and $L$-functions of meromorphic modular forms on
%  $\mathrm{SL}_2(\mathbb{Z})$ via finite-contour cocycles},
%  \emph{in preparation}.
  
\end{thebibliography}


%  ============================  SKELETON (final organization)  ============================
%  §1 Setup (prose) --- modularity = weight-k slash; S_k < M_k < M^!_k < F_k and S_k < S^!_k < M^!_k; n := k-2; q = e(tau), e(x) := e^{2 pi i x}.
%  §2 Motivation (prose) --- L-functions of weak cusp forms (BGKO; the KZ extra relation); BFK's deformed-Mellin regularization; the case for one uniform finite construction.  (respectful BFK framing)
%
%  §3 The L-integral   [sec:Lint]
%     def:segments    S- and T-segments at tau_0 (pole-avoiding)
%     def:kernel      Hurwitz T-kernel
%     def:Lint        L-integral L*(f,s; gamma^S, gamma^T) on F_k
%     def:Lfun        L-function L = (2 pi)^s / Gamma(s) * L*
%     lem:finite      finiteness on F_k
%     prop:indep      basepoint/path independence on M^!_k
%     rem:mellin      classical Mellin recovery (tau_0 -> i infty, cusp forms)
%
%  §4 The L-function   [sec:Lfun]
%     prop:HGG     Hurwitz-Gamma identity
%     thm:weakL       L* at tau_0 = i is the incomplete-Gamma sum; functional equation
%     def:cocycle     period cocycle C^f_gamma
%     lem:gate        parabolic representatives exist iff c_f(0) = 0
%     def:rf          period polynomial r_f = the parabolic S-value
%     prop:periods    period coefficients are special L-values
%     cor:quasi       ... and are Brown's quasi-periods (k = 12,16,18,20,22,26)
%     rem:legendre    period/quasi-period determinant lies in Q^x (2 pi i)^{k-1}
%     rem:canonical   [subsection] why a finite contour, not the Mellin transform
%
%  §5 Meromorphic forms   [sec:mero]   (poles split at the edge Im(tau) = 1)
%     def:proj        polylog projection (removes poles with Im(tau_p) >= 1)
%     def:cRf         regularized cusp coefficients c_{Rf}(n)
%     def:Finfty      the T-periodic space F(infty)
%     lem:commute     q-series commutes with the integral iff no pole has Im >= 1
%     def:polylog     polylog integral I_N(s,z)
%     lem:polylogcf   closed form of I_N for Im z > 1 (polylog inversion)
%     thm:mero        the explicit L-function formula (cusp sum + interior polylog integrals)
%     lem:onedge      a pole on the edge: principal value (half-residue; tau = i is the S-segment PV)
%     cor:wall        wall-crossing: L* is per-chamber constant in tau_0, jumping by the residue crossed
%
%  NUMERICS [VerifyScripts]:  verify_periods_all_weights.py (cor:quasi, rem:legendre);
%     check_inversion.py (lem:polylogcf);  check_offedge.py (thm:mero);
%     check_onedge.py + check_lambda.py (lem:onedge, cor:wall).
%  =====================================================================================


\end{document}

\begin{thebibliography}{99}

\bibitem{Brown2017}
F.\ Brown, \emph{A class of non-holomorphic modular forms III},
arXiv:1710.07912 (2017).

\bibitem{BFK2011}
K.\ Bringmann, K.-H.\ Fricke, and Z.\ Kent,
\emph{Special $L$-values and periods of weakly holomorphic modular forms},
Ramanujan J.\ \textbf{24} (2011), 119--139.

\bibitem{1806}
D.\ A.\ McGady,
\emph{$L$-functions for meromorphic modular forms and sum rules in
conformal field theory},
arXiv:1806.09874 (2018).

\bibitem{KZ1984}
W.\ Kohnen and D.\ Zagier,
\emph{Modular forms with rational periods},
in \emph{Modular Forms} (R.~A.\ Rankin, ed.), Ellis Horwood, 1984, pp.\ 197--249.

\bibitem{BCD2018}
R.\ Bruggeman, Y.\ Choie, and N.\ Diamantis,
\emph{Holomorphic automorphic forms and cohomology},
Mem.\ Amer.\ Math.\ Soc.\ \textbf{253} (2018), no.\ 1212; arXiv:1404.6718.

\bibitem{PasolPopa2013}
V.\ Pa\c{s}ol and A.\ A.\ Popa,
\emph{Modular forms and period polynomials},
Proc.\ London Math.\ Soc.\ (3) \textbf{107} (2013), 713--743; arXiv:1202.5802.

\bibitem{WangZhang2022}
W.\ Wang and H.\ Zhang,
\emph{Meromorphic quasi-modular forms and their $L$-functions},
arXiv:2202.09542 (2022).

\bibitem{BrownFonseca2025}
F.\ Brown and T.\ J.\ Fonseca,
\emph{Single-valued periods of meromorphic modular forms and a motivic
interpretation of the Gross--Zagier conjecture},
arXiv:2508.04844 (2025).

\bibitem{BDE2016}
K.\ Bringmann, N.\ Diamantis, and S.\ Ehlen,
\emph{Regularized inner products and errors of modularity},
Int.\ Math.\ Res.\ Not.\ IMRN (2016); arXiv:1603.03056.

\bibitem{DiamantisDrewitt2020}
N.\ Diamantis and J.\ Drewitt,
\emph{Period functions associated to real-analytic modular forms},
Res.\ Math.\ Sci.\ \textbf{7} (2020), no.\ 21.

\bibitem{VerifyScripts}
\emph{Numerical verification scripts} (companion code), 2026.

\end{thebibliography}

\end{document}
